Synthesizing knowledge from scientific literature is essential for uncovering new research directions, refining methodologies, and supporting evidence-based decisions. However, the vast volume of papers published annually makes it increasingly difficult for researchers to stay informed. Effective synthesis requires precise retrieval, accurate attribution, and real-time access to current literature. 
While large language models (LLMs) show promise in assisting researchers, they face significant challenges, including hallucinations~\citep{mallen2022not,mishra2024finegrained}, reliance on outdated pre-training data~\citep{kasai2022realtime}, and a lack of transparent attribution. For instance, when tasked with citing up-to-date literature, GPT-4 fabricated citations in 78-90\% of cases across fields like computer science and biomedicine in our experiments. 

Retrieval-augmented LMs~\citep{lewis2020retrieval,guu2020retrieval}, on the other hand, can mitigate many of these issues by integrating retrieved external knowledge sources at inference-time, driving the development of systems for literature search and synthesis ~\citep{agarwal2024litllm,zheng2024openresearcher,skarlinski2024language}. 
However, many such systems rely on black-box APIs or general-purpose LLMs that are neither optimized for literature synthesis nor paired with open, domain-specific retrieval datastores (i.e., a processed corpus and corresponding retrieval index) that are specifically suited for scientific domains. 
Moreover, evaluations for scientific literature synthesis have been limited, using single-discipline and small-scale human evaluations~\citep{agarwal2024litllm,zheng2024openresearcher} or simplified tasks such as multiple-choice question answering~\citep{skarlinski2024language}. 

To address these gaps, we present \model (Figure~\ref{fig:overview}, top), a state-of-the-art retrieval-augmented LM with a specialized paper datastore and retrievers trained on scientific literature. At inference time, \model retrieves relevant passages and uses iterative self-feedback generation to refine its own output. 
We further train a new, efficient 8B LM. 
To evaluate the effectiveness of \model, we introduce \data (Figure~\ref{fig:overview}, middle), a benchmark specifically designed to enable a realistic and reproducible evaluation of open-ended scientific question answering.


\model (Section~\ref{sec:model}) uses our new \model-\textsc{DataStore} (OSDS), which contains 45 million open-access papers from Semantic Scholar, along with 237 million corresponding passage embeddings. 
To the best of our knowledge, this is the largest open-sourced datastore of scientific domains. 
\model first retrieves passages from OSDS using a retriever and reranker. Subsequently, an LM synthesizes the retrieved passages to generate responses with citations.
\model iteratively refines its outputs through natural language feedback, which improves quality and adaptively incorporates supplementary information. 
This pipeline is also used to create large-scale, high-quality training data for smaller, more efficient models. We generate synthetic queries and instructions from sampled datastore passages, feed them into \model, and use intermediate and final output to train open 8B model, \model-8B and retrieval models. 

\data (Section~\ref{sec:dataset}) is a benchmark designed to evaluate the ability of models to understand and synthesize existing research. Unlike previous benchmarks that assume that answers can be found in a single paper (e.g., scientific fact-checking; \citealp{wadden-etal-2020-fact,skarlinski2024language}), many real-world scenarios require identifying multiple relevant papers and generating long-form output with accurate citations. To address these challenges, we curated a dataset of 2,967 literature synthesis questions, along with 208 long-form responses that are written by experts and span four scientific disciplines, namely computer science, physics, biomedicine, and neuroscience. These responses were crafted by Ph.D. students and postdoctoral researchers with more than three years of experience and relevant publications in the field. On average, each response required approximately one hour to compose.  
We also introduce a multifaceted evaluation protocol that combines automated metrics and human assessments to measure citation accuracy, factual correctness, content coverage, coherence, and overall quality. 
This multifaceted approach ensures robust and reproducible evaluations, both automatic and human-driven. 

We evaluated proprietary and open models (e.g., GPT4o, Llama 3.1 8B, 70B) with and without retrieval capabilities, as well as specialized systems like PaperQA2~\citep{skarlinski2024language}, on \data (Section~\ref{sec:results}). While GPT4o demonstrated strong general performance, it struggled with citation accuracy and coverage, often producing inaccurate or non-existent citations.
\model outperformed both LM-only and retrieval-augmented pipelines, surpassing proprietary and open-source systems. Notably, using fully open-source checkpoints, \model outperformed PaperQA2~\citep{skarlinski2024language}, built on proprietary LMs, and production systems like Perplexity Pro, achieving 6\% and 10\% improvements, respectively. Additionally, \model's use of smaller, efficient retrievers significantly reduced costs.
Combining \model with GPT4o also improved correctness by 12\% over GPT4o alone.
The \model pipeline can also enhance off-the-shelf LMs. For example, when using GPT-4o as the underlying model, \model-GPT4o achieves a 12\% improvement in correctness compared to GPT-4o alone. 
 
In addition to automatic evaluations on \data, we conducted detailed expert assessments with 16 scientists from fields such as computer science, physics, and biomedicine (Section~\ref{sec:human_eval}). These experts performed pairwise and fine-grained evaluations of \model's outputs against 108 expert-written responses to literature synthesis queries in \data. 
\model, when paired with GPT-4o and our trained 8B model, consistently outperformed expert-written responses, with win rates of 70\% and 51\%, respectively. 
In contrast, GPT-4o without retrieval struggled with information coverage and was rated as less helpful than human experts, achieving only a 31\% win rate against human responses.  
This highlights that \model-generated outputs are more comprehensive, well-organized, and useful for synthesizing literature.
These findings demonstrate that \model produces high-quality outputs that are not only competitive with expert-written answers but, in some cases, exceed them, particularly in terms of coverage and organization.

\model-8B is an open retrieval-augmented LM that avoids reliance on proprietary LMs or retrieval systems, leveraging one of the largest datastores in scientific literature domains. 
We release the full \model ecosystem, including code, trained retrievers, the LM checkpoint, the datastore, the \data benchmark, expert evaluation tools, and a public demo.
